\chapter{Introduction}
\label{ch:intro}

\section{Problem domain}

GraphHunter is an interactive threat-hunting tool built around a
single idea: most attacker behaviour is naturally a path through a
directed temporal graph, so let the analyst \emph{ask for paths}
rather than write per-table SQL. A hypothesis is expressed in a
small arrow-chain DSL --- for example,

\begin{center}
\code{User -[Auth \{status="Failure"\}]-> IP HAVING count(distinct IP)
>= 3 WITHIN 24h}
\end{center}

and the engine returns the paths in an ingested graph that satisfy
it, ranked by composite anomaly score. The graph itself is built
from heterogeneous log sources (Microsoft Sentinel, Defender XDR,
Fortinet, Sysmon, IIS, generic CSV/JSON). Typical workloads range
from a few thousand edges (one incident ticket) to hundreds of
millions (multi-week EVTX dump from a medium enterprise).

\section{System overview}

The codebase is a Rust monorepo with five top-level components,
layered as in Figure~\ref{fig:crate-graph}.

\begin{figure}[h]
\centering
\begin{tikzpicture}[
  node distance=4mm and 10mm,
  crate/.style={draw, rounded corners, minimum width=3.2cm,
                minimum height=0.85cm, align=center, font=\small},
  layer/.style={draw, dashed, rounded corners, inner sep=4mm},
]

  \node[crate] (core) {\textbf{graph\_hunter\_core}\\ engine};
  \node[crate, right=of core] (api) {\textbf{graph\_hunter\_api}\\ canonical façade};

  \node[crate, above right=of api] (tauri) {Tauri desktop\\ (\texttt{app/src-tauri})};
  \node[crate, right=of api]       (http)  {HTTP server\\ (\texttt{http\_api.rs})};
  \node[crate, below right=of api] (mcp)   {MCP server\\ (\texttt{graph-hunter-mcp})};

  \node[crate, below=of core] (libgm) {\textbf{libgraphmatch}\\ C++ AVX2 SIMD};

  \draw[-{Stealth}] (core) -- (api);
  \draw[-{Stealth}] (api) -- (tauri);
  \draw[-{Stealth}] (api) -- (http);
  \draw[-{Stealth}] (api) -- (mcp);
  \draw[-{Stealth}] (libgm) -- (core) node[midway,right,font=\tiny]{FFI};
  \draw[-{Stealth}, dashed] (mcp) to[bend right=25] node[above,font=\tiny]{HTTP} (http);

\end{tikzpicture}
\caption{Crate graph. The engine (\texttt{graph\_hunter\_core}) is
consumed by a canonical API façade (\texttt{graph\_hunter\_api});
three transport layers delegate to that façade. The MCP server
reaches the same engine by calling the HTTP server over the
loopback interface rather than linking to the façade directly.}
\label{fig:crate-graph}
\end{figure}

The engine is a pure Rust library: no framework dependencies, no
global state beyond a handful of \texttt{OnceLock}-cached catalogs.
The façade crate (\code{graph\_hunter\_api}) wraps the engine with
session management, DTOs, and an \code{EventEmitter} trait, so the
three transport layers can share a single implementation without
importing the low-level types themselves. The C++ SIMD backend
(\code{libgraphmatch}) is linked via FFI when the \code{simd}
feature is enabled; a pure-Rust \code{std::arch} alternative
(\code{simd\_rust}) is shipping as its replacement (Chapter~\ref{ch:simd}).

\section{What this reference covers}

The document walks the engine bottom-up:

\begin{itemize}
    \item \textbf{Part II} --- the data structures the engine stores
        in memory: compact relations, the bidirectional string
        interner, the in-memory graph, the spillable edge store,
        the external metadata store, and the hypothesis AST the DSL
        parser produces.
    \item \textbf{Part III} --- the core algorithms: DSL parsing,
        temporal pattern matching (DFS with pruning, smart
        anomaly-bounded variant), SIMD acceleration, anomaly
        scoring, centrality analytics (Brandes betweenness, temporal
        PageRank), path deduplication \& pagination, and hunt diff.
    \item \textbf{Part IV} --- the streaming ingest pipeline:
        \code{LogSource} trait, parsers, the bounded-channel writer,
        MVCC edge-level stamping, and overflow persistence.
    \item \textbf{Part V} --- the extension seams introduced by the
        P2 refactor: \code{GraphRead}/\code{GraphWrite},
        \code{ScoreComponent}/\code{Scorer},
        \code{StringInternerBackend}, \code{IngestAdapter}, and the
        YAML catalog loader.
    \item \textbf{Part VI} --- the transport layer: the canonical
        API, DTOs, \code{EventEmitter}, Tauri/HTTP/MCP adapters, the
        parity harness, and session persistence.
\end{itemize}

Four appendices close the document: a one-page complexity summary
(App.~\ref{app:complexity}), the ATT\&CK catalog reference
(App.~\ref{app:catalog}), a benchmark table
(App.~\ref{app:bench}), and a source index mapping every algorithm
to its \code{file:line} (App.~\ref{app:source}).

\begin{inthecode}
Crate roots are \filepath{graph\_hunter\_core/Cargo.toml},
\filepath{graph\_hunter\_api/Cargo.toml},
\filepath{app/src-tauri/Cargo.toml},
\filepath{graph-hunter-mcp/package.json}, and
\filepath{libgraphmatch/CMakeLists.txt}.
\end{inthecode}
